\documentclass[11pt,letterpaper]{article}
\usepackage[margin=0.82in]{geometry}
\usepackage[T1]{fontenc}
\usepackage{lmodern,microtype}
\usepackage{amsmath,amssymb,graphicx,booktabs,array,float}
\usepackage[font=small,labelfont=bf]{caption}
\usepackage[hidelinks]{hyperref}
\usepackage{url}
\usepackage{placeins}
\setlength{\parindent}{0pt}
\setlength{\parskip}{6pt}
\setlength{\textfloatsep}{14pt}
\setlength{\intextsep}{12pt}
\setlength{\emergencystretch}{2em}
\raggedbottom
\renewcommand{\arraystretch}{1.15}
\graphicspath{{figures/}{code/}{../../outputs/figures/}{../assets/final/}}
\newcommand{\codeimage}[2]{\begin{figure}[H]\centering\includegraphics[width=\linewidth]{#1}\caption*{Executed Python excerpt: #2. Corresponding results follow below.}\end{figure}}
\newcommand{\resulttable}[2]{\begin{table}[H]\centering\small\input{generated/#1.tex}\caption{#2}\end{table}}
\newcommand{\plot}[4][0.96]{\begin{figure}[H]\centering\includegraphics[width=#1\linewidth]{#2}\caption{#3}\label{#4}\end{figure}}
\hypersetup{pdftitle={Robust Detection of Anomalous Taxi Fares Using Least Trimmed Squares Regression},pdfauthor={Aditya Kumar}}
\begin{document}
\begin{center}
{\Large\bfseries Robust Detection of Anomalous Taxi Fares\\Using Least Trimmed Squares Regression\par}
\vspace{7pt}
Aditya Kumar \quad IED/10007/23\\
ED409 Regression Techniques \quad Instructor: Dr. Tina Dutta
\end{center}

\section*{Business case and dataset}
Suppose a taxi audit team receives thousands of trip records each day. Its task is to identify journeys whose recorded fares need a closer look. A large fare by itself is not enough to justify a flag: it may simply belong to a long journey. The more useful question is whether the fare makes sense alongside the recorded distance, duration, and other trip details.

There is a second difficulty. The records used to learn a normal fare pattern can themselves contain unusual observations. For example, a journey recorded as lasting almost a day may have a fare of only a few dollars. If a regression allows such records to determine its coefficients, the resulting benchmark can become less representative of ordinary trips. The team then risks comparing new transactions against a distorted standard.

This case examines Least Trimmed Squares (LTS) regression as a way to obtain a more stable fare benchmark. Ordinary least squares (OLS) provides the comparison. Both methods use the same information, but LTS initially fits the observations most compatible with a common relationship. The business purpose is to understand that relationship, assess its stability, and produce a list of records for manual review.

\subsection*{What the dataset contains}
The study uses the \emph{January 2026 Yellow Taxi Trip Records} published by the New York City Taxi and Limousine Commission (TLC). It contains 3,724,889 submitted trip records. One row represents one recorded journey. The records are supplied by authorized technology providers; TLC does not guarantee their accuracy (NYC TLC, n.d.).

The response variable is \texttt{fare\_amount}, the metered time-and-distance fare in US dollars. It is different from the passenger's total bill: tips, tolls, and separately itemized extras are recorded in other fields. The explanatory variables are trip distance in miles, elapsed duration in minutes, passenger count, pickup time, and weekend status. Duration is calculated from the pickup and drop-off timestamps.

After basic eligibility checks, 2,278,082 records remain. A reproducible sample of 50,000 trips is divided into 40,000 training observations and 10,000 evaluation observations. No upper limit is imposed on fare, distance, duration, or passenger count. Keeping large values is essential because their effect on regression is the subject of the case.

\small
\textbf{Official source:} \url{https://www.nyc.gov/site/tlc/about/tlc-trip-record-data.page}\\
\textbf{File:} \url{https://d37ci6vzurychx.cloudfront.net/trip-data/yellow_tripdata_2026-01.parquet}
\normalsize

\newpage
\section*{Business problems to be addressed}
The audit team needs answers to four questions. Together, they establish what the fare benchmark describes, how trustworthy it is, and how its warnings should be used. The dataset does not contain verified fraud labels or independently established correct fares. A model discrepancy must therefore be treated as a reason to investigate, rather than evidence of misconduct.

\subsection*{1. What is the fare relationship followed by ordinary trips?}
How does metered fare vary with distance and duration after allowing for passenger count, pickup time, and weekend status? The analyst needs an interpretable benchmark, with coefficients expressed in familiar units such as dollars per mile and dollars per minute. Fitting OLS and LTS to the same predictors will show whether they describe a similar relationship or whether unusual observations substantially change that description. These coefficients are statistical associations; they are not estimates of the official tariff or causal effects of extending a trip.

\subsection*{2. How much do unusual trips affect OLS?}
Do extreme distance--duration--fare combinations pull the OLS fit away from the pattern followed by most journeys? This question matters because the audit team wants a useful comparison for ordinary transactions, but it must also understand how the model handles rare ones. The analysis will examine residuals and influential observations, then compare OLS and LTS using both absolute-error and squared-error measures. A method can fit typical trips more closely while making larger errors on a few extremes, so one performance number cannot settle this question.

\subsection*{3. Which records should enter an audit queue?}
Which training observations disagree strongly with the resistant fare relationship, and how large would the resulting review workload be? The analysis will use LTS residuals and reweighting decisions to identify audit candidates. The queue should include fares that are unusually high or unusually low relative to their trip characteristics. Before drawing conclusions, an investigator would need to check the timestamps, meter record, route, and fare context. The purpose is to prioritize that work, not to automate accusations.

\subsection*{4. Does LTS remain stable when some fares are corrupted?}
What happens if exactly 5\% of the training fares are increased by \$100? This controlled experiment provides a direct check of resistance. It selects 2,000 training rows using seed 42, changes only their fare values, and refits both methods. The original evaluation fares remain the reference for measuring prediction performance. Because the shifted rows are known, the experiment can also count how many are detected and how many unaltered records are flagged alongside them.

The case supports a decision about a stable fare benchmark and an explainable audit queue. Its conclusions apply to eligible January 2026 standard-rate trips and the specified corruption experiment; it does not estimate a fraud rate.

\newpage
\section{Preparing a consistent set of trips}
Before comparing estimators, the observations must describe a reasonably consistent business process. Standard metered trips are therefore analysed separately from other rate codes. A stable \texttt{source\_row\_id} preserves the connection between each selected trip and its position in the original file.

The filters are applied in sequence. The removal count at each step consequently refers only to records that passed the earlier steps. The executed excerpt below loads the fixed modelling split and reads the preparation audit already produced by the data-preparation workflow.

\codeimage{code_provenance.png}{loading the modelling split and preparation audit}
\resulttable{filter_audit}{Sequential eligibility checks on the original file.}

The January condition fixes the reporting period, and rate code 1 restricts the fare regime. Positive distance and fare define the positive-fare journey model used here. Duration must be at least one minute, and passenger count must be at least one. A missing value in a required filter field fails the relevant rule. These decisions exclude refunds, zero-distance records, and incomplete observations from the present case; they do not establish that those records are necessarily erroneous.

There is deliberately no upper-bound filter. In particular, the long-duration observations remain available for the regression comparison. Removing them because they affect OLS would decide the robustness question before fitting LTS.

\subsection{From the eligible population to the modelling sample}
A seed-42 permutation selects 50,000 of the 2,278,082 eligible records. A second fixed seed-42 procedure assigns 40,000 rows to training and 10,000 to evaluation. The earlier 10,000-row preparation sample is a separate exploratory artifact; it is not the sample behind the final results.

\subsection{What each predictor means}
Distance and duration describe the completed journey. Passenger count adds occupancy information. Pickup time enters as a pair of cyclical variables,
\[
s_i=\sin(2\pi h_i/24),\qquad c_i=\cos(2\pi h_i/24),
\]
where $h_i$ is the pickup hour. This preserves the fact that 23:00 and 00:00 are close on a clock. A weekend indicator equals one for Saturday or Sunday. The identification fields and split label are not predictors.

The model uses information known at trip completion. Its business role is retrospective fare review; it is not a model for quoting a fare at pickup.

\section{What the training data reveal}
The 40,000 training observations have complete values for the exported analytical fields. Table~\ref{tab:descriptive} shows that most trips are short and inexpensive, but the distributions extend far beyond their medians.
\begin{table}[H]\centering\small\begin{tabular}{lrrr}
\toprule
Variable & Median & 99th percentile & Maximum \\
\midrule
Fare amount USD & 12.10 & 65.30 & 204.60 \\
Trip distance miles & 1.53 & 15.50 & 54.83 \\
Duration minutes & 11.58 & 50.38 & 1,437.05 \\
Passenger count & 1.00 & 4.00 & 6.00 \\
\bottomrule
\end{tabular}
\caption{Distribution summaries for the 40,000 training observations.}\label{tab:descriptive}\end{table}

The duration distribution is particularly revealing. Its median is about 11.58 minutes, yet one record lasts 1,437.05 minutes. A regression fitted to all squared errors can respond strongly to such a large predictor value. The record may reflect an unusual operation, an interrupted meter, or a timestamp problem; the available fields cannot determine which explanation is correct.

\plot{training_distributions.png}{Training distributions. The histograms show values through the 99.5th percentile for readability; full maxima are printed on the axes. All observations remain in the analysis.}{fig:distributions}

The relationship plots help explain why duration and distance should be considered together. Fare generally rises with both, but a journey can take longer without travelling much farther. The broadening of the fare--duration cloud also suggests that one constant residual variance will be inadequate.

\plot{training_fare_relationships.png}{Fare relationships in the central training data. Darker hexagons contain more trips. The visual limits do not remove observations from any fitted model.}{fig:relationships}

Pearson fare--distance correlation is 0.953. For fare and duration, Pearson correlation is only 0.494, while Spearman rank correlation is 0.967. Pearson uses the actual magnitudes; Spearman uses their ranks. The gap is consistent with a strong monotonic pattern among ordinary trips alongside a few extremely long durations that do not carry correspondingly large fares. The two measures describe different aspects of the data, and neither implies causation.

\FloatBarrier
\section{How OLS and LTS fit the same fare model}
Both estimators use an intercept and the same six predictors. In original units, the common specification is
\begin{align*}
y_i={}&\beta_0+\beta_1 d_i+\beta_2 t_i+\beta_3 p_i\\
      &+\beta_4 s_i+\beta_5 c_i+\beta_6 w_i+\varepsilon_i,
\end{align*}
where $y_i$ is metered fare, $d_i$ is distance, $t_i$ is duration, $p_i$ is passenger count, and $w_i$ denotes a weekend. This common specification makes the comparison about the estimator's response to the observations.

\subsection{Why trimming can change the fitted relationship}
For a trial coefficient vector $\beta$, let $r_i(\beta)=y_i-x_i^\top\beta$. OLS chooses the vector that minimizes all squared residuals:
\[
\widehat\beta_{\mathrm{OLS}}=\arg\min_\beta\sum_{i=1}^{n}r_i^2(\beta).
\]
A residual twice as large contributes four times as much to this objective. Large response errors, combined with unusual predictor values, can therefore alter the fitted coefficients substantially.

LTS first orders the squared residuals for each trial fit and minimizes only the smallest $h$:
\[
\widehat\beta_{\mathrm{LTS}}=\arg\min_\beta\sum_{i=1}^{h}r_{(i)}^2(\beta),
\qquad r_{(1)}^2\leq\cdots\leq r_{(n)}^2.
\]
The retained observations are the ones most compatible with that fit; they are not the trips with the smallest fares. The subset can change as the trial coefficients change (Rousseeuw \& Van Driessen, 2006).

The primary setting is $\alpha=0.75$, so $h=30,000$ of the 40,000 training rows enter the trimmed objective. This does not assert that the other 25\% are bad records. It controls how much of the data must support a common relationship. Prespecified sensitivity settings are 0.50 and 0.90. At 0.50 the software enforces its minimum $h=20,003$; at 0.90, $h=36,000$.

\subsection{How the computation is performed}
Searching every possible subset would be impractical. The project uses FAST-LTS through RobPy version 0.0.6, with 500 random starting subsets and seed 42. For each promising fit, a concentration step selects the $h$ smallest squared residuals and refits least squares to those observations. Repeating the step reduces the trimmed objective. The best result among the starts is retained. This efficient search approximates LTS for the large dataset; a global optimum is not guaranteed (RobPy, n.d.).

All predictor means and standard deviations are estimated from training rows. The same transformation is applied to evaluation predictors. Fare remains in dollars. The following executed code fits the primary models and reports the raw subset size and robust residual scale.

\codeimage{code_fit.png}{training-only scaling and the primary fit}
\begin{table}[H]\centering\begin{tabular}{lr}\toprule
Output & Value\\\midrule
Raw LTS subset $h$ & 30,000\\
MAD residual scale (USD) & 0.510606\\
Observations retained for reweighting & 33,548\\\bottomrule
\end{tabular}\caption{Output of the primary clean fit.}\end{table}

\subsection{Why there is a second fit after trimming}
Raw LTS supplies the resistant starting relationship. The project then uses its residuals to identify observations close enough to that relationship for a final least-squares refit. This step is called reweighting. It can use more observations than the raw trimmed subset and improve efficiency.

Here the scale is the normal-consistent median absolute deviation (MAD),
\[
\widehat s=\frac{\operatorname{median}_i|r_i-\operatorname{median}(r)|}{0.67449}.
\]
The standardized residual is $z_i=r_i/\widehat s$. A row receives weight one if $|z_i|\leq2.241403$ and zero otherwise. The cutoff is $\sqrt{\chi^2_{1,0.975}}$. Least squares is refitted using the weight-one observations. This is the project's explicit MAD-based rule: RobPy's built-in reweighting is disabled, and no claim of exact finite-sample calibration is made.

\codeimage{code_lts_algorithm.png}{raw FAST-LTS, MAD scaling, and the explicit refit}
The executed result is again $h=30,000$, $\widehat s=0.510606$, and 33,548 positive weights. The assertion checks that the explicit implementation reproduces the shared fitting helper. The raw subset size and the final reweighting count have different roles: the first defines the trimmed objective, and the second determines the final refit. Both leave all 40,000 records available for error measurement and audit screening.

When a model is fitted to standardized predictors, its coefficients are converted back before interpretation. If $b_j^*$ is a standardized slope, the original-unit slope is $b_j=b_j^*/s_j$. The original intercept is $b_0=b_0^*-\sum_j b_j^*\bar x_j/s_j$, where $\bar x_j$ and $s_j$ are training quantities.

\section{What the OLS diagnostics tell us}
The training fit is also used to check whether conventional OLS uncertainty is credible. The code below fits OLS and computes HC3 standard errors, a leverage-adjusted member of the White heteroskedasticity-consistent family taught in Module I.
\codeimage{code_diagnostics.png}{OLS and HC3 uncertainty}
\resulttable{hc3}{OLS coefficients and standard errors. HC3 changes the uncertainty estimate, not the coefficients.}

For distance, the standard error increases from about 0.006 under the constant-variance formula to 0.096 with HC3. For duration it increases from about 0.001 to 0.039. The coefficients do not change; their estimated precision does. This warns against treating the ordinary OLS standard errors as reliable in these data.

\plot{training_residual_diagnostics.png}{Training OLS residual diagnostics. The left plot shows changing residual spread; the normal Q--Q plot on the right shows substantial tail departures.}{fig:residuals}

The plots give the practical reason for that warning. Residual variability increases with fitted fare, and the Q--Q plot departs sharply from a normal reference. Residual skewness is $-3.21$ and ordinary kurtosis is 187.40, compared with normal kurtosis 3. Residual RMSE rises from \$1.91 in the lowest fitted-fare decile to \$6.22 in the highest.

Supporting tests agree with these patterns. The Breusch--Pagan LM statistic is 33,831.46, with $p<10^{-16}$, against the null of constant variance. Ramsey RESET gives $F=3,692.56$, also with $p<10^{-16}$, indicating that the additive mean form is incomplete. These are named implementations of assumption checking, rather than additional central methods. With 40,000 observations, small $p$-values alone would not establish a practically important problem; the plots and error magnitudes provide that context.

\plot{training_binned_residuals.png}{Average OLS residuals in predictor quantile bins. Systematic departures from zero indicate that the additive linear mean leaves structure unexplained. The bands are approximate 95\% intervals for bin means.}{fig:binned}

Figure~\ref{fig:binned} makes the mean-form issue easier to see. Distance residuals change from negative to positive and then back again. Duration residuals rise from about $-\$2.48$ to $+\$4.64$. A correctly represented conditional mean should not leave such a persistent pattern. The shared linear specification remains useful for an interpretable estimator comparison, but it is an approximation to the complete fare process.

Predictor VIF values range from 1.01 to 1.21. VIF measures how much linear overlap among predictors inflates coefficient variance; values this close to one suggest little inflation. Influence is a more visible concern. Cook's distance screens how much an observation can change the OLS fit by combining its residual with its predictor leverage. There are 1,663 training rows above the screening value $4/n$. The four highest values, about 47.77--56.62, belong to records lasting roughly 1,412--1,437 minutes with fares between \$7.20 and \$24.00. Such configurations can strongly affect the OLS duration slope. They remain in the dataset because their correctness cannot be settled from the available fields.

HC3 does not fix the mean function or prove independence of trips. Likewise, Durbin--Watson near two in randomly ordered records cannot establish absence of temporal dependence. These limits are carried into the conclusions.

\FloatBarrier
\section{Clean-data results: two different notions of a good fit}
The first comparison uses genuine recorded fares, before the artificial stress test. Here, ``clean'' means that no artificial fare shift has been added; it does not mean that the original records are verified error-free. The comparison asks how closely each estimator fits ordinary journeys and how it handles the complete distribution, including the retained extremes. The code predicts on the fixed 10,000-row evaluation sample after model fitting.
\codeimage{code_clean.png}{prediction on the evaluation sample}
\resulttable{clean_performance}{Clean-model errors. Dollar units apply to MAE, median absolute error, and RMSE. LTS denotes the primary reweighted fit at $\alpha=0.75$.}

\subsection{Interpreting the error measures}
For evaluation errors $e_i=y_i-\widehat y_i$, MAE averages $|e_i|$, whereas median absolute error takes their median. RMSE is $\sqrt{\operatorname{mean}(e_i^2)}$, so a few very large misses can dominate it. Evaluation $R^2$ compares squared prediction error with the squared deviations from the evaluation response mean:
\[
R^2=1-\frac{\sum_i(y_i-\widehat y_i)^2}{\sum_i(y_i-\bar y_{\mathrm{evaluation}})^2}.
\]

LTS has much smaller absolute errors. Its evaluation MAE is \$1.015, compared with \$2.086 for OLS; median absolute error is \$0.347, compared with \$1.594. It therefore describes the fare pattern followed by many ordinary trips more closely.

The ranking reverses when large errors receive a squared penalty. OLS has RMSE \$3.253 and $R^2=0.923$, while LTS has RMSE \$8.756 and $R^2=0.445$. LTS resists unusual training configurations and can consequently make large errors on similar extreme evaluation records. OLS accommodates those records more strongly, reducing squared loss while moving away from the central pattern. Figure~\ref{fig:cleanmetrics} displays both sides of this trade-off.

\plot{clean_model_performance.png}{Evaluation MAE and RMSE for OLS and the raw and reweighted LTS sensitivity fits. Lower bars indicate smaller errors. The opposite MAE and RMSE rankings are part of the result.}{fig:cleanmetrics}

\subsection{What changed in the estimated fare relationship}
\resulttable{coefficients}{Original-unit coefficients for the clean primary fits.}
OLS estimates about \$3.622 per mile and \$0.082 per minute, holding the other included predictors fixed. Reweighted LTS estimates about \$1.813 per mile and \$0.605 per minute. Thus the two estimators allocate the distance and duration contributions very differently. The influential duration records help explain why the OLS time slope is small.

The small passenger and time-cycle coefficients under LTS indicate a much weaker contribution within its fitted relationship. They are not supported here by conventional OLS significance tests on the selected subset. None of the slopes should be interpreted as an official meter charge: actual metering depends on operating conditions that this additive specification does not fully represent.

These findings answer the first two business questions. LTS describes a different dominant fare relationship and reduces typical absolute errors; OLS achieves the better overall squared-error compromise on this heterogeneous sample. A business user must decide which loss and which population matter for the intended decision.

\subsection{Sensitivity to the retained fraction}
\resulttable{sensitivity}{Evaluation results for all prespecified LTS settings.}
At $\alpha=0.90$, reweighted MAE is \$0.848 and RMSE \$8.280; at $\alpha=0.50$, they are \$1.043 and \$8.825. Retaining more observations changes the fitted compromise. The primary value remains 0.75 because it was fixed before comparison. Selecting the best alpha from the evaluation table would turn that sample into tuning data. Raw and reweighted errors are close at each alpha, showing that the principal difference from OLS comes from the resistant fit rather than a large gain from the final refit.

\section{Turning LTS residuals into an audit queue}
The third business question concerns review decisions, rather than prediction accuracy alone. A trip becomes an audit candidate if its raw-LTS reweighting weight is zero or its absolute standardized robust residual exceeds 2.5.

The two conditions are nested. The zero-weight threshold is about 2.241, so every residual above 2.5 already has zero weight. They do not provide two independent confirmations of an anomaly. In the clean training data, 6,452 rows have zero weight and 5,929 exceed 2.5. Their union is therefore 6,452 records, or 16.13\% of training trips.

That is a substantial queue. The saved candidate table ranks records by absolute standardized residual and retains their source identifiers, timestamps, fare, and trip attributes. A sensible first action would be to inspect a small ranked sample, checking timestamp consistency and meter context before deciding whether the screen is useful operationally. The data do not tell us how many flags represent a true recording error, an omitted trip characteristic, or a legitimate unusual journey.

The screen therefore answers \emph{which records disagree with the benchmark}. It cannot answer \emph{which records prove fraud}. The next experiment supplies a known artificial label so that detection performance can be evaluated under a specified form of corruption.

\FloatBarrier
\section{Stress test: adding \$100 to 5 percent of training fares}
Exactly 2,000 training rows are selected without replacement using seed 42. Their responses are changed according to
\[
y_i^*=y_i+100\,I_i,
\]
where $I_i=1$ for a selected row and zero otherwise. Predictor values and training scaling remain fixed. Both models are refitted and compared against the same original evaluation fares. This is response contamination: it does not simulate corrupted distances or timestamps.
\codeimage{code_stress.png}{constructing the shift, checking its scope, and refitting}
\resulttable{stress_performance}{Clean evaluation performance before and after corrupting training fares.}

OLS deteriorates substantially. Evaluation MAE rises from \$2.086 to \$5.383, RMSE rises from \$3.253 to \$5.939, and $R^2$ falls from 0.923 to 0.745. Its coefficients have absorbed part of the artificial inflation.

Reweighted LTS remains close to its original performance. MAE changes from \$1.015 to \$1.000 and RMSE from \$8.756 to \$8.737. These small improvements reflect the particular selected observations and changed fit; they do not show that corruption is beneficial. The main observation is the small amount of change compared with OLS.

\plot{contamination_performance.png}{Clean evaluation errors before and after response contamination. The nearly equal LTS bars show stability; OLS errors increase markedly.}{fig:stressmetrics}

\subsection{Why the OLS fit moves upward}
\resulttable{coefficient_change}{Changes in key original-unit coefficients. Change means the stressed coefficient minus the clean coefficient.}
The average response increase in training is $0.05\times\$100=\$5$. The OLS intercept rises by \$4.827, close to that shift. Its distance and time-related coefficients also move because the selected rows are not perfectly balanced across every predictor.

The LTS intercept changes by only $-\$0.025$. Its distance slope changes by about \$0.020 per mile and its duration slope by about $-\$0.000542$ per minute. This provides a direct answer to the fourth business question: under the specified corruption, LTS preserves its fitted fare relationship much more closely.

\plot{contamination_coefficient_change.png}{Absolute changes in slope coefficients. The intercept is omitted from this plot and reported separately in the table. Coefficients have different original units, so bar heights should be compared between estimators within a term, not used to rank predictor importance.}{fig:coefficientchanges}

\subsection{Seeing the observations behind the comparison}
The next two figures separate the location of the artificial points from the movement of the models. This makes the ordinary fare pattern readable without compressing it underneath the much larger shifted fares.

\plot[0.94]{actual_contamination_overview.png}{The artificial fare band in actual comparable training trips. Grey dots are unaltered observations; red dots are selected trips after exactly \$100 is added. The modelling experiment shifts 2,000 of all 40,000 training rows. Only the comparable display group is shown here.}{fig:outlierband}

For the display, trips are restricted to one passenger, weekdays, pickup hours 11--13, and distance between 1.18 and 1.88 miles. This group contains 1,038 trips. The picture samples 250 unaltered trips using seed 42 and includes all 48 selected corruptions in that group. These restrictions affect only the picture. All models were fitted to the full 40,000 training rows.

The fitted lines evaluate the original six-predictor equations at 1.53 miles, one passenger, weekday, and noon while duration varies. The observed points are close to these conditions. These are conditional fitted lines, not new regressions fitted to the display group.

\plot[1.0]{actual_points_clean_stress_fitted_lines.png}{The ordinary fare range before and after the shift, using identical axes. Grey dots are observed fares. Each panel shows only two lines: OLS in blue and reweighted LTS in orange. The right-hand arrow marks the OLS prediction increase at a fixed duration; the shifted red observations are above this zoom and are shown in Figure~\ref{fig:outlierband}.}{fig:lines}

\begin{samepage}
In Figure~\ref{fig:lines}, the orange LTS line follows the visible duration--fare pattern in both panels. The blue OLS line is relatively flat before corruption and shifts upward after it. Showing one OLS and one LTS line in each panel makes the change easier to compare than overlaying four almost coincident curves.

The picture illustrates one part of the data. Its clear local pattern does not replace the full-sample error measures, which also account for the retained extremes.
\end{samepage}

\subsection{What the detection results establish}
\codeimage{code_detection.png}{counting flags against known artificial labels}
\begin{table}[H]\centering\begin{tabular}{lr}\toprule
Experiment outcome & Count or rate\\\midrule
Shifted records flagged (true positives) & 2,000\\
Shifted records missed (false negatives) & 0\\
Unaltered records flagged (false positives) & 5,684\\
Unaltered records unflagged (true negatives) & 32,316\\
Recall & 100.00\%\\
Precision & 26.03\%\\
False-positive rate & 14.96\%\\\bottomrule
\end{tabular}\caption{Detection against the known artificial-contamination label.}\end{table}

Recall is the proportion of shifted rows detected: $\mathrm{TP}/(\mathrm{TP}+\mathrm{FN})$. Precision is the proportion of flags that are shifted rows: $\mathrm{TP}/(\mathrm{TP}+\mathrm{FP})$. LTS finds every artificial corruption, but it also flags many unaltered observations. Hence recall is perfect for this large shift while precision is only 26.03\%.

A ``false positive'' here means unaltered in the experiment. It does not establish that the row is an ordinary, correctly recorded journey. Some of these records may already be unusual relative to the fitted benchmark. This distinction explains why the experiment demonstrates resistance and screening sensitivity without validating automatic fraud classification.

\clearpage
\section{What the audit team should conclude}
The case began with the need for a dependable comparison for ordinary taxi fares. LTS supplies a useful description of the dominant relationship: its absolute errors are smaller, and its distance--duration coefficients remain stable when a minority of training responses are shifted by a large amount. This supports its use as a descriptive benchmark and as one input to an audit queue.

OLS retains an important advantage in this dataset. Its clean and stressed RMSE values are lower than the corresponding LTS values. It accommodates genuine extreme records more strongly, producing a better overall squared-error compromise. Resistance to corruption and minimum squared prediction loss are therefore distinct goals. The business recommendation depends on which goal matters for the proposed use.

For an audit team, the next practical step would be a limited manual review of ranked candidates, with confirmed outcomes recorded for later threshold calibration. The 16.13\% clean flag rate is too large to treat casually as a list of established errors. The present analysis cannot estimate the financial value of an audit programme or the prevalence of misconduct.

\subsection{Limits that matter for interpretation}
The evidence comes from one month, one rate-code population, and one fixed random split. Trips may share traffic conditions, routes, drivers, or recording systems, so random row splitting does not establish independent errors or future-month accuracy.

The additive fare model leaves visible mean structure and changing residual variance unexplained. Robust fitting does not repair every aspect of model specification. The fixed MAD rule is an implementation choice whose audit workload depends on this heterogeneous population. Conventional OLS inference is not transferred to the LTS coefficients after subset selection.

Finally, the stress test is one deterministic realization of a conspicuous upward response shift. It includes no repeated-run uncertainty interval and no predictor corruption. Its strong LTS recall should not be generalized to subtle errors, downward shifts, or coordinated changes in trip attributes. Within those boundaries, the experiment gives clear evidence that LTS resists the specified response contamination better than OLS.

\clearpage
\section{Direct answers to the business questions}
The following answers bring together the evidence for the four questions posed at the start. Both estimators were fitted to 40,000 training trips; prediction errors refer to the same 10,000 evaluation trips.

\subsection{What is the fare relationship followed by ordinary trips?}
The visible pattern is that fare rises with distance and duration. Holding the other included predictors fixed, primary reweighted LTS estimates \$1.813 per additional mile and \$0.605 per additional minute. OLS estimates \$3.622 per mile and \$0.082 per minute, allocating much more of the relationship to distance. LTS follows the central trip pattern more closely: its median absolute prediction error is \$0.347, compared with \$1.594 for OLS. These are conditional statistical associations within the chosen model, not official tariff rates or causal effects.

\subsection{How much do unusual trips affect OLS?}
The diagnostics identify substantial influence: 1,663 training observations exceed the Cook's-distance screen, and the four most influential records combine roughly day-long durations with fares of only \$7.20--\$24.00. They help explain the relatively flat OLS duration relationship. LTS reduces evaluation MAE from the OLS value of \$2.086 to \$1.015. However, OLS has lower RMSE: \$3.253 versus \$8.756. Thus OLS accommodates extreme configurations more strongly, while LTS fits many ordinary trips more closely. The evidence supports different strengths, rather than one universally better estimator.

\subsection{Which records should enter an audit queue?}
The screen identifies 6,452 training trips, or 16.13\%, that have zero reweighting weight or an absolute standardized raw-LTS residual above 2.5. The stricter 2.5 rule flags 5,929 of these trips; it does not add separate candidates to the union. Review should begin with the largest absolute standardized deviations, covering fares unusually high or low for their recorded trip characteristics. Investigators should check timestamps, meter records, routes, and fare context. These are audit candidates; the available data cannot establish which flags are recording errors or fraud.

\subsection{Does LTS remain stable when some fares are corrupted?}
Yes, under the specified experiment. Adding \$100 to exactly 2,000 training fares raises OLS evaluation MAE from \$2.086 to \$5.383 and shifts its intercept upward by \$4.827. LTS MAE changes only from \$1.015 to \$1.000, with an intercept change of about $-\$0.025$. Its audit screen detects all 2,000 shifted rows, but also flags 5,684 unaltered rows, giving 100\% recall and 26.03\% precision against the artificial labels. LTS therefore demonstrates stability and sensitivity to this large response shift; it does not establish precise fraud identification or resistance to every type of corruption.

\clearpage
\section*{References}
\begingroup
\raggedright
\everypar{\hangindent=1.5em\hangafter=1}
New York City Taxi and Limousine Commission. (n.d.). \emph{TLC trip record data} [Data set]. Retrieved November 1, 2026, from \url{https://www.nyc.gov/site/tlc/about/tlc-trip-record-data.page}

New York City Taxi and Limousine Commission. (2025, March 18). \emph{Data dictionary: Yellow taxi trip records}. \url{https://www.nyc.gov/assets/tlc/downloads/pdf/data_dictionary_trip_records_yellow.pdf}

Rousseeuw, P. J., \& Van Driessen, K. (2006). Computing LTS regression for large data sets. \emph{Data Mining and Knowledge Discovery, 12}, 29--45. \url{https://doi.org/10.1007/s10618-005-0024-4}

RobPy. (n.d.). \emph{robpy.regression} [Software documentation]. \url{https://robpy.readthedocs.io/en/stable/api/robpy.regression.html}
\par\endgroup

\end{document}
